\documentclass[twocolumn]{aastex631}
\begin{document}
\title{A residual reionization ionizing-photon-budget shortfall for star-forming galaxies at z\~{}10 (beta systematic corner)}
\author{NebulaMind Lab (autonomous pipeline)}
\affiliation{NebulaMind Lab, \url{https://lab.nebulamind.net}}
\begin{abstract}
Reionization ionizing-photon-budget reconciliation at z\~{}10: star-forming galaxies require f\_esc=0.782 (+0.788/-0.400) to close the budget (Madau-Dickinson SFRD, log xi\_ion=25.5+/-0.15, clumping C=2-5, JWST-SFRD tail), versus indirect-proxy-inferred f\_esc=0.050 (+0.076/-0.030) from LzLCS O32/beta calibrations. Median delta(required-inferred)=+0.708 dex-frac (16-84\%: +0.303 to +1.498); 99\% of the systematic MC shows a shortfall. Conclusion: a genuine SHORTFALL remains, and the sign holds under both O32 and beta calibrations. The result is bounded by the xi\_ion x clumping x proxy-calibration systematic, not by statistics. Generated autonomously from public data (jwst) via the NebulaMind Lab runner: a bounded, reproducible, descriptive study.
\end{abstract}
\section{Introduction}
Recent studies have highlighted a potential challenge in our understanding of the reionization process, suggesting that star-forming galaxies may struggle to produce sufficient ionizing photons to account for the observed reionization [Muñoz2024]. This issue has been explored through various approaches, including assessments of the galaxy ionizing photon budget at z < 10 [Duncan2015] and excursion set reionization models [Park2022]. However, a more nuanced reconciliation of the reionization-photon-budget is necessary to address this discrepancy.
\section{Data and method}
To investigate this problem, we employed a literature-anchored budget calculation that did not rely on any survey catalog data. Instead, we used the cosmic SFRD from Madau \& Dickinson (2014) as one possible model and adopted published values for xi\_ion and the O32/beta f\_esc proxy calibrations [LzLCS: Chisholm+22, Flury+22; Simmonds+24]. Our method focused on calculating the ionizing-photon-budget to determine if star-forming galaxies can account for reionization under this specific SFRD model.
\section{Result}
Our analysis revealed that at z\~{}10, star-forming galaxies require an escape fraction of f\_esc=0.782 (+0.788/-0.400) to close the budget when considering the Madau-Dickinson SFRD, log xi\_ion=25.5+/-0.15, and clumping factor C=2-5. However, indirect-proxy-inferred escape fractions from LzLCS O32/beta calibrations yield a significantly lower value of f\_esc=0.050 (+0.076/-0.030). This discrepancy results in a median delta(required-inferred)=+0.708 dex-frac (16-84\%: +0.303 to +1.498), with 99\% of systematic Monte Carlo simulations showing a shortfall.
\begin{figure}\centering\includegraphics[width=\columnwidth]{result.png}\caption{A residual reionization ionizing-photon-budget shortfall for star-forming galaxies at z\~{}10 (beta systematic corner)}\end{figure}
\section{Caveats}
It is essential to acknowledge the limitations and uncertainties associated with our approach. Our analysis relies on the Madau-Dickinson SFRD model, which may not capture the full range of possible star formation histories. Additionally, we recognize that our adopted xi\_ion values and clumping factor are subject to uncertainties, and a more robust error propagation for these parameters is needed in future work. Furthermore, our analysis does not account for other sources of ionizing photons, such as active galactic nuclei or X-ray binaries, which could contribute to the overall photon budget. A more comprehensive understanding of these factors, including alternative SFRD models and their implications, is necessary to resolve the reionization-photon-budget crisis.
\section*{References}
\footnotesize
[Muoz2024] Muñoz et al. (2024). Reionization after JWST: a photon budget crisis?. bibcode 2024MNRAS.535L..37M \par [Davies2021] Davies et al. (2021). The Predicament of Absorption-dominated Reionization: Increased Demands on Ionizing Sources. bibcode 2021ApJ...918L..35D \par [Duncan2015] Duncan et al. (2015). Powering reionization: assessing the galaxy ionizing photon budget at z < 10. bibcode 2015MNRAS.451.2030D \par [Park2022] Park et al. (2022). Calibrating excursion set reionization models to approximately conserve ionizing photons. bibcode 2022MNRAS.517..192P \par [Madau2017] Madau (2017). Cosmic Reionization after Planck and before JWST: An Analytic Approach. bibcode 2017ApJ...851...50M

\end{document}
